% Generated draft: requires mathematical and source review.
\documentclass[11pt]{article}
\usepackage[margin=25mm]{geometry}
\usepackage{fontspec}
\setmainfont{Latin Modern Roman}
\usepackage{amsmath,amssymb,mathtools,unicode-math}
\setmathfont{Latin Modern Math}
\usepackage{enumitem,needspace,xurl,hyperref}
\setlength{\parindent}{0pt}
\setlength{\parskip}{5pt}
\setlength{\emergencystretch}{3em}
\newcommand{\N}{\mathbb{N}}
\newcommand{\Z}{\mathbb{Z}}
\newcommand{\Q}{\mathbb{Q}}
\newcommand{\R}{\mathbb{R}}
\newcommand{\C}{\mathbb{C}}
\newcommand{\F}{\mathbb{F}}
\newcommand{\A}{\mathbb{A}}
\newcommand{\PP}{\mathbb{P}}
\newcommand{\E}{\mathbb{E}}
\newcommand{\eps}{\varepsilon}
\newcommand{\dd}{\,\mathrm{d}}
\begin{document}
\section*{UnsolvedMath ID: 11100026}
\textbf{Source problem code:} catalogue entry
\subsection{Definitions and mathematical statement}
For a prime $p$ and height $n\geq1$, Morava $K$-theory conventionally has coefficients $K(n)_*\cong\F_p[v_n,v_n^{-1}]$ with $|v_n|=2(p^n-1)$. Hovey's Applications item 2, prompted by Neil Strickland, asks for Morava $K$-theory of moduli spaces appearing in differential geometry and physics, beyond their rational homology. It names gauge theory and also mentions solitons and Seiberg--Witten theory.

The source explicitly does not identify the gauge-theoretic moduli space precisely. A common framework starts with connections $\mathcal A(P)$ on a bundle $P$ and a gauge group $\mathcal G(P)$; a moduli problem may impose differential equations and then take a quotient. But no base manifold, structure group, bundle, equations, regularity or boundary conditions are supplied, so this framework is context rather than a definition of the intended space.

The research request is to compute the appropriate $K(n)_*(\mathcal M)$ or $K(n)^*(\mathcal M)$ for specified such moduli problems, with prime, height and homology/cohomology conventions made explicit. Coarse quotients, homotopy quotients or stacks, compactifications and compactly supported theories can give different objects and cannot be selected silently. The source asks an exploratory question, not a universal vanishing or nontriviality assertion.

The registry reports scattered ordinary-$K$-theory and chromatic work but no systematic answer or directly matching 2024--2026 computation located by its auditor. Because the target spaces are unspecified, that search assessment is not a theorem that no relevant Morava-$K$ computations exist.
\subsection{Short English statement}
Explore Morava $K$-theory of gauge-theoretic, soliton and Seiberg--Witten moduli spaces, after specifying the actual moduli problem and quotient. The source deliberately leaves its gauge-theory example unidentified.
\subsection{Sources}
\begin{itemize}
\item Local registry: \url{data/unsolvedmath/problems.json}, permanent UnsolvedMath ID $11100026$, exact source code \texttt{catalogue entry}, attributed to Mark Hovey with the suggestion credited to Neil Strickland.
\item \emph{Hovey's Algebraic Topology Problem List}, Applications item 2: \url{https://www-users.cse.umn.edu/~tlawson/hovey/}. Supplied access date: 2026-07-29; extraction: html-ordered-list; 65 top-level items in nine category pages.
\item No specific moduli-space construction, primary computation or bibliographic reference is supplied. The registry mentions instanton and ordinary-$K$-theory examples only descriptively and reports an unsuccessful 2024--2026 matching-result search.
\end{itemize}
\end{document}
